\documentclass[11pt]{article}
\usepackage[margin=1in]{geometry}
\usepackage{amsmath,amssymb,amsthm}
\usepackage{booktabs}
\usepackage{microtype}
\usepackage[hidelinks]{hyperref}
\usepackage{enumitem}
\usepackage{array}

\newtheorem{assumption}{Assumption}
\newtheorem{proposition}{Proposition}
\newtheorem{remark}{Remark}
\newcommand{\E}{\mathbb{E}}
\newcommand{\R}{\mathbb{R}}
\newcommand{\ind}{\mathbb{1}}
\newcommand{\henv}{h_{\mathrm{env}}}
\newcommand{\doop}[1]{\mathrm{do}(#1)}

\title{\textbf{The Portability CIN:\\ Calibrated Counterfactual Projection of NHL Player Impact\\ under Environment Interventions}}
\author{OptBot Research \\ \small Technical Report v1.0 --- July 2026 --- Confidential}
\date{}

\begin{document}
\maketitle

\begin{abstract}
\noindent We present the Portability Causal Intervention Network (CIN), a system that projects
a professional hockey player's on-ice impact in an environment he has never played in ---
a new team, line, and role --- before a single shift is observed there. The central design
thesis is that credible causal AI is achieved not by architectural scale but by four
enforceable commitments: (C1) answering only intervention classes that occur as natural
experiments in the data, with tiered refusal outside support; (C2) structural deconfounding
of player talent from deployment via a frozen, out-of-fold talent channel; (C3) optimizing
the model on the distribution it will be queried on, by up-weighting windows adjacent to
observed context switches; and (C4) conformal calibration of uncertainty on the intervention
distribution itself, with achieved coverage attached to every emitted interval. On a
frozen-clock backtest of 635 real player--team transitions (2021--2025), the minimal
instantiation of the system (v0: destination-environment expectation plus a shrunken
residual talent prior) attains RMSE 5.30 on first-40-game on-ice expected-goal share,
versus 5.80 for a Marcel-style industry baseline ($-8.6\%$; bootstrap 95\% CI on the
difference $[-0.67, -0.33]$, excluding zero) and 7.40 for last-season carryover. Emitted
80\% intervals achieve 79.8\% empirical coverage. We give the complete data substrate,
estimators, neural architecture, leakage-audit methodology (including three data corruptions
whose discovery materially changed the system), and an 18-item threat model.
\end{abstract}

\section{Introduction}

Player acquisition in the National Hockey League is a capital-allocation problem executed
under a specific epistemic handicap: all available statistics describe a player's output
\emph{in his previous environment}. On-ice results in hockey are violently contextual ---
linemates, opposition, deployment, coaching system, and schedule jointly determine a large
fraction of observed variance --- so the acquisition decision is, implicitly, a causal
question: \emph{what would this player's output be under our environment?} Formally it is a
counterfactual query, and the industry answers it today with heuristics (scouting intuition)
or context-free extrapolations (weighted averages of past seasons).

This report documents a system that answers the query directly and, more importantly,
\emph{auditable-ly}. Our position is that ``causal AI'' fails in practice for reasons that
are disciplinary rather than architectural: systems answer interventions their data never
contained, entangle treatment assignment with subject quality, extrapolate without support
diagnostics, and report uncertainty that has never been checked against realized frequencies.
Sections~\ref{sec:formulation}--\ref{sec:serving} present a design in which each of these
failure modes is closed by an explicit, testable mechanism rather than by optimism.

Two results anchor the report. First, the minimal version of the system already clears the
strongest public baseline with a confidence interval excluding zero
(Section~\ref{sec:eval}). Second, the engineering audit that preceded any modeling
(Section~\ref{sec:data}) discovered three silent data corruptions --- including an outcomes
column that was identically zero across 761{,}823 rows --- whose repair retroactively
explains the failure of a previous uncertainty model. We consider the audit methodology as
much a contribution as the architecture.

\section{Problem Formulation}\label{sec:formulation}

\subsection{Units and notation}

The atomic observational unit is the \emph{window}: a maximal interval of game time with
constant on-ice personnel and manpower state, indexed $w$. Each window carries, for a focal
player $p$ present in it:
\begin{itemize}[nosep]
\item exposure $s_w > 0$ (the focal player's seconds on ice in $w$);
\item an environment description $e_w = (c_w,\, T_w,\, V_w)$: pre-faceoff context $c_w$
(score state, zone start, period, rink, schedule, manpower, team-level priors), the teammate
multiset $T_w$ with co-presence seconds, and the opponent multiset $V_w$;
\item a lever vector $\ell_w \in \R^{17}$ of coach-controlled deployment quantities;
\item an outcome vector $Y_w \in \R^{10}$: on-ice expected goals for/against (xGF, xGA),
goals, shots, and four microstatistics, all attributed over $s_w$.
\end{itemize}
Windows aggregate to player-games and player-seasons in the obvious exposure-weighted way.
We write $Y^{60}$ for per-60-minute rates and
$\mathrm{xGF}\% = \mathrm{xGF}^{60} / (\mathrm{xGF}^{60} + \mathrm{xGA}^{60})$.

\subsection{The estimand}

Fix a player $p$ with history $\mathcal{H}_p(t_0)$ up to a freeze date $t_0$, and a
\emph{hypothetical} environment distribution $\mathcal{E}'$ --- in the product, the empirical
window distribution of a destination team's line slot, with teammates optionally pinned by
the user. The target is the interventional mean over a horizon $H$ (first 40 games):
\begin{equation}
\theta(p, \mathcal{E}', t_0) \;=\;
\E\!\left[\, \bar{Y}^{60}_{p,H} \;\middle|\; \doop{E \sim \mathcal{E}'},\; \mathcal{H}_p(t_0) \right],
\label{eq:estimand}
\end{equation}
reported as $\mathrm{xGF}\%$ with an interval $[\theta_L, \theta_U]$ whose nominal coverage
must match realized coverage on historical interventions (Section~\ref{sec:conformal}).

\subsection{Identification}

Equation~\eqref{eq:estimand} is identified from observational data under the following
assumptions, each of which is engineered toward rather than assumed:

\begin{assumption}[Overlap on the intervention class]\label{as:overlap}
$\mathcal{E}'$ lies within the support of environments realized by players comparable to
$p$. \emph{Enforcement:} the intervention class is restricted to environment swaps ---
of which the ledger contains 729 natural instances --- and a support gate refuses queries
outside observed ranges (Section~\ref{sec:support}).
\end{assumption}

\begin{assumption}[Conditional ignorability given talent and context]\label{as:ignor}
$Y_w \perp E_w \mid (z_p, c_w, \ell_w)$, where $z_p$ is player ability. Deployment is
assigned by coaches \emph{on their beliefs about ability}; conditioning on an explicit,
externally estimated ability quantity $z_p$ is what renders the assignment ignorable.
\emph{Enforcement:} $z_p$ is estimated out-of-fold from a baseline that never observes
player identity (Section~\ref{sec:talent}) and enters the model through a dedicated channel
(C2), so environment coefficients cannot absorb it.
\end{assumption}

\begin{assumption}[Consistency / no interference across the swap]\label{as:consistency}
The player's ability transports across environments over the projection horizon; general
equilibrium effects of the acquisition on the rest of the league are negligible at
window scale.
\end{assumption}

\begin{remark}
Assumption~\ref{as:ignor} is the load-bearing one and it is only as good as $z_p$. The
residual confounding from imperfect $z_p$ is not argued away: it is \emph{priced} --- the
conformal layer is fit on realized intervention errors and therefore absorbs, empirically,
whatever ignorability violations survive (Section~\ref{sec:threats}, T1--T3).
\end{remark}

\section{Relation to Prior Approaches}

\textbf{Marcel-class projections} (weighted multi-season averages with regression to the
mean and an age adjustment) are the de facto industry standard; they are environment-blind
by construction and serve as our primary baseline. \textbf{RAPM/WAR-family models} estimate
context-adjusted player value by ridge regression over on-ice indicator matrices; they share
our goal of separating player from context but produce retrospective value, not forward
projections under specified interventions, and are afflicted by the collinearity of
persistent line combinations (our T3). \textbf{Deep sequence models} of play-by-play data
optimize predictive likelihood and inherit the deployment confound wholesale. Within the
causal-ML taxonomy, our system is closest to an S-learner with an explicitly engineered
adjustment set, wrapped in split-conformal inference; the novelty we claim is not any single
component but the enforcement architecture binding them (C1--C4), which we have not seen
deployed jointly in a production sports system.

\section{Data Substrate and Forensic Repairs}\label{sec:data}

\subsection{Perfect Windows: a contracted training surface}

All modeling reads from a single table, \texttt{perfect\_windows\_v2} (6{,}490{,}872 rows;
seasons 2018-19 through 2024-25; 5v5), assembled by joining audited sources and validated at
build time against five machine-checked contracts:
\begin{enumerate}[nosep]
\item every player-performance feature is lagged and empirical-Bayes shrunk, and ships with
its effective sample size $n_{\mathrm{eff}}$;
\item every context feature is measurable at the window's opening faceoff (PRE-timing);
\item people lists (\texttt{with\_ids}, \texttt{vs\_ids}) are fixed length $K{=}5$, sorted
by co-presence seconds descending (a legacy system silently consumed \emph{unsorted} lists
as if they were top-$K$; the raw sources are confirmed unsorted);
\item baseline predictions are strictly out-of-fold, with fold provenance retained;
\item outcome columns are targets and are refused as features by the validator.
\end{enumerate}

\subsection{Forensic findings}

Three corruptions were found by the validation gate before any training run, in data that
had already been used by predecessor systems:
\begin{enumerate}[nosep]
\item \textbf{Null outcomes.} The player-game observation file underlying all residual
computation carried $y \equiv 0$ in every one of 761{,}823 rows --- a silently failed join.
Every ``residual'' was $-\mu$. This retroactively explains the predecessor uncertainty
model's 12\% empirical coverage at a nominal 80\% (its $\bar z = -1.87$ is exactly the
missing-actuals artifact). Outcomes were rebuilt from per-game ground truth with a 100.0\%
join rate across seven seasons.
\item \textbf{Duplicate exposure.} The 10.29M-row window spine contained 1{,}491{,}039
full-row duplicates spanning all eight seasons (a combine-step defect), double-counting
exposure downstream; deduplication left zero key conflicts. The corruption had propagated
into the out-of-fold prediction files ($12.9\%$ duplicate keys), which are now key-deduped
at every load.
\item \textbf{Global baseline miscalibration.} The environment baseline underpredicts
on-ice xG by a season-dependent factor (mean residual $+0.9$ to $+1.9$ per 60). Rather than
rescale the baseline (a label-aware operation with leakage risk), talent is defined
\emph{league-relatively} by per-season demeaning (Section~\ref{sec:talent}), rendering the
estimator invariant to any global miscalibration.
\end{enumerate}

\section{The Talent Prior}\label{sec:talent}

\subsection{Construction}

Let $r_{pg}$ denote player $p$'s game-$g$ on-ice xGF$^{60}$ residual against the
out-of-fold, identity-blind environment baseline, demeaned within season $s(g)$ by the
TOI-weighted league mean:
\begin{equation}
\tilde r_{pg} = r_{pg} - \frac{\sum_{q,h:\, s(h)=s(g)} w_{qh}\, r_{qh}}{\sum w_{qh}},
\qquad w_{pg} = \text{TOI}_{pg}\ \text{(hours)} .
\end{equation}
Talent state ahead of any date $t$ is an exponentially decayed, \emph{lagged} running
moment set with half-life $\lambda$ (default 40 games):
\begin{equation}
m_p(t) = \frac{\sum_{g: t_g < t} \rho^{\,\Delta(g,t)} w_{pg}\, \tilde r_{pg}}
              {\sum_{g: t_g < t} \rho^{\,\Delta(g,t)} w_{pg}},
\qquad
n_p(t) = \sum_{g: t_g < t} \rho^{\,\Delta(g,t)} w_{pg},
\qquad \rho = 2^{-1/\lambda},
\end{equation}
implemented so the state recorded at each game excludes that game (verified by unit test:
the running-statistic recorder emits before ingesting). The served quantity is
empirical-Bayes shrunk toward the league (zero, post-demeaning):
\begin{equation}
z_p(t) = \frac{n_p(t)}{n_p(t) + K}\; m_p(t),
\qquad
\widehat{\mathrm{se}}_p(t) = \max\!\left(\sqrt{v_p(t)/n_p(t)},\ \mathrm{se}_{\min}\right),
\label{eq:shrink}
\end{equation}
with $K$ selected by out-of-time predictive risk over a grid (never hand-set) and a hard
floor $\mathrm{se}_{\min}$ replacing a legacy pathology in which 89.6\% of low-$n$ rows
carried a collapsed, identically zero standard error. Defensive talent uses $-\,$xGA
residuals with the same machinery.

\subsection{Face validity}

On the first correctly computed residuals in the project's history, the untuned prior ranks
McDavid first, Draisaitl second, and MacKinnon third among players with
$n_{\mathrm{eff}} > 15$ TOI-hours, with Kucherov, Crosby (age 37), Makar, and Barzal in the
top fifteen and Kucherov's \emph{defensive} component negative. We treat this as a
necessary sanity property, not a result; the results are Section~\ref{sec:eval}.

\section{The Environment Tower}\label{sec:tower}

\subsection{Design principle: the focal player has no embedding}

The predecessor encoder assigned every player a learned identity vector and then fought a
counterintelligence war (gradient-reversal scrubbers, stream firewalls) to keep that vector
from memorizing deployment and team. We remove the disease instead of treating it: the
focal player enters the network only through $\big(z_p, \widehat{\mathrm{se}}_p,
n_p\big)$ and explicit bio features (age, age$^2$, handedness, position). There is no
lookup keyed on his identity, hence nothing to scrub. Teammates and opponents \emph{do}
receive embeddings --- from the focal player's vantage point, other players are environment,
and environment is precisely what the intervention swaps.

\subsection{Architecture}

Inputs are restricted to a frozen legal schema (Section~\ref{sec:legal}). Categorical
context $c_w$ maps through embeddings; continuous team priors pass through a normalized
linear block. The people payload is encoded by seconds-biased attention pooling: for
teammate embeddings $u_1,\dots,u_K$ with co-presence seconds $s_1 \ge \dots \ge s_K$,
\begin{equation}
\alpha_k = \mathrm{softmax}_k\!\big( u_k^\top q + \log \max(s_k, 1) \big) \cdot
\ind\{s_k > 0\},
\qquad
h_{\mathrm{with}} = \textstyle\sum_k \alpha_k u_k,
\label{eq:setpool}
\end{equation}
and symmetrically for opponents. The tower output is
$\henv = f_\theta\big(c_w, h_{\mathrm{with}}, h_{\mathrm{vs}}\big) \in \R^{256}$.
Two training-time corruptions are integral to the design: \emph{teammate dropout}
(random masking of WITH identities) prevents $\henv$ from fingerprinting the focal player
through his habitual linemates and simultaneously matches the serving regime, where lineups
are partially hypothesized; unseen identities at serving time resolve to pooled
position$\times$role fallback vectors, with degradation under synthetic unknowns measured
as an explicit test.

\subsection{The legal input schema}\label{sec:legal}

Every candidate feature was classified by a single criterion --- \emph{can the scenario
builder reproduce it for a hypothetical destination from pre-$t_0$ information?} --- into
\textsc{synth} (exactly constructible), \textsc{marg} (served as the destination slot's
empirical distribution via bootstrap), or \textsc{banned}. Notable rulings, each traced to
the generating code rather than inferred from names:
\begin{itemize}[nosep]
\item \texttt{matchup\_quality\_pct} --- \textsc{banned}. Traced construction: opponents'
\emph{full-game} TOI percentile within role (includes post-window deployment, i.e.\ coach
reactions to the game's own outcomes), weighted by \emph{realized within-window}
head-to-head overlap seconds. Matchup difficulty reaches the tower legally through opponent
identities and lagged opponent priors instead.
\item \texttt{standing\_prior} --- \textsc{banned}. One code path selects standings rows
with \texttt{after\_pk $\le$ current gamePk}, admitting the game's own result.
\item \texttt{ai\_OZ\_start}, \texttt{ai\_DZ\_start} --- \textsc{allowed}. ``ai'' denotes
\emph{after icing}: an icing-forced zone start, fully determined at the window's first
second, and incidentally the cleanest deployment natural experiment hockey offers.
\item All in-window shot-style counts and sums --- \textsc{banned} as outcomes wearing
feature costumes, notwithstanding their (real) predictive value: the scenario builder can
never supply them, so any accuracy they contribute is train/serve fiction.
\end{itemize}
The rulings compile to an \texttt{assert\_legal()} guard that refuses banned features at
trainer construction, citing the trace in the raised error.

\subsection{Acceptance probes}\label{sec:probes}

A checkpoint ships only if a five-probe scorecard passes, all probes fit on held-out splits
(the dry-run on random embeddings exposed in-sample probe overfitting --- $R^2 = 0.095$ on
pure noise --- which would have flunked clean encoders; probes were repaired before first
use and validated to refuse a random encoder, catch a synthetically leaky one, and pass a
synthetically clean one):
\begin{center}
\begin{tabular}{@{}llll@{}}
\toprule
Probe & Readout from & Requirement & Rationale \\
\midrule
P1 & $\henv \to$ team identity & accuracy $< 3\times$ chance & env must not memorize employer \\
P2 & $\henv \to$ focal OZ share & $R^2 < 0.10$ & env must not encode focal deployment \\
P3 & $\henv$/talent $\to$ next-season ability & $R^2 > 0.30$ & the system retains real signal \\
P4 & cross-season retrieval & top-5 $> 0.5$ & representations stable across years \\
P5 & twin-distance ratio & $< 1.5$ & player moves must not fracture identity \\
\bottomrule
\end{tabular}
\end{center}

\section{The CIN Core}\label{sec:core}

\subsection{Trunk}

With $x = \big[\henv;\ \phi_z(z_p, \widehat{\mathrm{se}}_p, n_p);\ \phi_b(\mathrm{bio})\big]$
projected to width $d{=}768$, the trunk applies $L{=}4$ pre-norm residual blocks with
SwiGLU activation, each modulated by the lever vector through feature-wise linear
modulation (FiLM):
\begin{equation}
\mathrm{Block}(h) = h + \mathrm{Drop}\Big(
\big(1 + \tanh W_g \ell\big) \odot \mathrm{SwiGLU}\big(\mathrm{LN}(h)\big) + W_b \ell \Big).
\end{equation}
FiLM gives levers clean intervention semantics --- knobs modulate the computation rather
than mixing additively with state --- and the talent channel $\phi_z$ is a separate pathway
whose contribution the trunk cannot re-derive from deployment inputs (C2 made structural).

\subsection{Heads and objective}

Ten heads emit $(\mu_j, \log\sigma_j)$ per window. The per-head loss is heteroscedastic
Gaussian NLL for the semi-continuous xG rates and Poisson-type likelihoods for count heads
(goals, shots, microstatistics), exposure-offset by $\log(s_w/60)$; 68.8\% of windows carry
$y_{\mathrm{xGF}} = 0$, and a homoscedastic Gaussian on that margin is misspecified.
Cross-task weights are learned (Kendall--Gal):
\begin{equation}
\mathcal{L} = \sum_{j=1}^{10} e^{-\eta_j}\,
\overline{\mathcal{L}_j}^{\,(\omega)} + \tfrac{1}{2}\eta_j ,
\qquad
\omega_w = \underbrace{s_w}_{\text{exposure}} \cdot
\underbrace{\big(1 + (\kappa - 1)\,\ind\{g_w \le 10\}\big)}_{\text{switch weight (C3)},\ \kappa = 3},
\end{equation}
where $g_w$ counts games since the focal player's last team change: the loss concentrates
on post-transition windows, the regime the product queries. Internal $\sigma_j$ are used
for training and uncertainty \emph{ranking} only; they are contractually barred from
human-facing intervals (Section~\ref{sec:conformal}). Training uses era-split checkpoints
(strictly-before 2021/22/23/24) so that every backtested projection is freeze-legal, and a
gradient-boosted twin trained on the identical legal features serves as a referee: the
neural core does not ship unless it beats both the twin and the v0 estimator under the
evaluation protocol of Section~\ref{sec:eval}.

\section{Serving}\label{sec:serving}

\subsection{Scenario construction and the v0 estimator}

Given $(p, t_0, \text{destination}, \text{line})$, the builder identifies the slot
incumbent by pre-$t_0$ TOI rank, bootstrap-resamples $N{=}2000$ of the incumbent's real
pre-$t_0$ windows (preserving the joint distribution of context --- score states, zone
starts, schedule --- rather than sampling marginals independently), overwrites the
traveling-player block with $p$'s frozen talent and bio, and pins user-specified linemates
into the WITH slots. The minimal estimator is then
\begin{equation}
\widehat{\mathrm{xGF}}^{60} = \E_{\hat{\mathcal{E}}'}\big[\mu^{\mathrm{OOF}}_{\mathrm{xGF}}\big] + z^{\mathrm{off}}_p(t_0),
\qquad
\widehat{\mathrm{xGA}}^{60} = \E_{\hat{\mathcal{E}}'}\big[\mu^{\mathrm{OOF}}_{\mathrm{xGA}}\big] - z^{\mathrm{def}}_p(t_0),
\label{eq:v0}
\end{equation}
i.e.\ destination environment expectation plus transported talent --- explainable in one
sentence and, per Section~\ref{sec:eval}, already sufficient to clear the industry baseline.
The neural path (v1) replaces the first terms with trunk forward passes over the same
synthetic windows, and is what makes linemate-level swaps and Tier-2 lever queries
answerable at all: estimator~\eqref{eq:v0} is structurally blind to linemate identity.

\subsection{Support gating}\label{sec:support}

Queries are tiered. Tier~1 (environment swap): answered, subject to slot support
($\ge 200$ real windows) and prior quality checks. Tier~2 (lever nudge): answered only
within the focal archetype's observed lever range $[q_{05}, q_{95}]$ harvested from the
validated lever-sweep corpus. Tier~3: the engine returns a refusal object,
\texttt{REFUSED\_OOS}, never a number. Refusal-as-output is, in our view, the single
clearest behavioral difference between causal inference as engineering and as branding.

\subsection{Conformal calibration}\label{sec:conformal}

Intervals come from split-conformal inference on the \emph{intervention population}: signed
errors of frozen-at-$t_0$ projections on the historical ledger, binned by talent
$n_{\mathrm{eff}}$, with the standard exchangeability-based marginal coverage guarantee
within bins. Every emitted interval carries its bin's \emph{achieved} coverage; emission
without it is a contract violation. Current global figures: nominal 80\%, achieved 79.8\%.
The predecessor system's internal-variance intervals achieved 12\% at the same nominal
level --- the difference between modeled and measured uncertainty is the credibility of the
entire product.

\section{Evaluation}\label{sec:eval}

\subsection{Protocol}

The ledger comprises every player--team transition 2021--2025 detected in our own data
(teamId change across consecutive appearances), qualified at $\ge 1000$ career 5v5 minutes
pre-move and $\ge 20$ GP post-move: 729 moves (196 in-season trades, 533 off-season).
For each move, the world is frozen at $t_0$: talent snapshot, destination templates, and
baselines all read strictly pre-$t_0$ data (era-split models for any learned component).
The target is realized on-ice xGF\% over the first 40 post-move games. An executable
leakage audit re-runs sampled moves with post-$t_0$ rows physically deleted and asserts
bit-identical projections. Baselines: last-season carryover; Marcel (5/4/3 season weights,
regression to league mean by minutes, age adjustment); and Marcel plus half the
team-quality delta (the strawman a hostile reviewer proposes).

\subsection{Results}

\begin{center}
\begin{tabular}{@{}lccc@{}}
\toprule
Projector & RMSE (xGF\%, first 40 GP) & $\Delta$ vs.\ Marcel & 95\% CI on $\Delta$ \\
\midrule
Last-season carryover & 7.47 & $+1.67$ & --- \\
Marcel & 5.80 & --- & --- \\
\textbf{CIN v0 (this work)} & \textbf{5.30} & $\mathbf{-0.50}$ ($-8.6\%$) & $[-0.67,\, -0.33]$ \\
\bottomrule
\end{tabular}
\end{center}

\noindent $n = 635$ scoreable moves; the confidence interval is a 10{,}000-resample
bootstrap of the paired RMSE difference and excludes zero by a wide margin. Sliced by move
type, v0 attains 5.14 on in-season trades versus 5.37 on off-season moves --- the system is
\emph{strongest} on the natural experiments closest to its identification story, and to the
customer's pain. Emitted 80\% intervals achieve 79.8\% coverage; interval half-widths
($\approx \pm 6.6$ points) are wide and honest, dominated by the sampling noise floor of a
40-game on-ice share ($\approx$3--4 points), which no projector can beat and which we
disclose rather than obscure. The v0 result is a floor in a second sense: the estimator
uses one scalar per player per side and an environment \emph{average}; the neural core
exists to convert linemate identity and lever structure into margin beyond $-8.6\%$, and
ships only if it does (Section~\ref{sec:core}).

\section{Threat Model}\label{sec:threats}

The complete 18-item audit (deployment confounding; talent--residual circularity; linemate
collinearity; ledger survivorship at the $\ge$20-GP filter, a collider whose bias applies
symmetrically to all compared projectors; acquisition selection-on-fit; injury
unobservables; role-assignment error; regression to the mean; era drift and short seasons;
goaltender contamination of defensive signal; score-state effects; out-of-support
extrapolation; the sampling noise floor; the three data corruptions of
Section~\ref{sec:data}; baseline integrity; deadline-context shift; cold-start prior
degeneracy; and uncertainty miscalibration) is maintained in the repository README with a
disposition for each: \textsc{solved} structurally, \textsc{priced} into conformal
intervals, \textsc{gated} at the API, or \textsc{open} with scheduled work. The
methodological stance is uniform: no threat is argued away; each is closed, measured, or
refused.

\section{Limitations and Ongoing Work}

(i) The scalar talent prior cannot express \emph{fit} --- complementarities between player
styles; v1 adds a low-dimensional, probe-gated style vector on the talent side, computed
from pre-$t_0$ windows only. (ii) The talent--outcome loop is one-step EM; iterating with
talent-aware refits is scheduled once the v1 harness exists. (iii) Shrinkage constant $K$
selected at the grid boundary under a per-game criterion; refit against next-season
aggregates is queued. (iv) Sensitivity of the headline to the survivorship filter
($\ge$10~GP variant) and to residual duplicate effects in legacy $\mu$ aggregation are
scheduled re-runs. (v) The league-translation layer for AHL/ECHL projection --- the
commercial extension --- is an exercise in data acquisition on this architecture, not new
methodology: the query class is identical.

\section*{Reproducibility}

The system is a single contracted repository (validation gate, executable schemas,
era-split training protocol, freeze-audited backtest, hash-sealed forward predictions at
\texttt{seal/}); every number in this report is regenerated by numbered pipeline scripts
\texttt{00}--\texttt{06} from raw sources, and forward projections for the 2026-27 season
are frozen and SHA-256--sealed before puck drop, to be scored publicly.

\end{document}
